% !PRINTED 596
\weberpage{596}

%<<CONT
sich auf die Einheitsgruppe reduciren. Es muss also jedenfalls
einmal der Fall eintreten, dass die Gruppe $P$ durch Adjunction
einer Wurzel einer reinen Gleichung reducirt wird.

Die Untersuchung dieser Frage wird ausserordentlich vereinfacht,
wenn man sie noch etwas umformt\,\footnote{Auf diese Form der Fragestellung hat zuerst C. \gesperrt{Jordan} hingewiesen (Trait\'e des substitutions p. 386).}.

Wir haben im \S.\,162 gesehen, dass die reinen Gleichungen zu den
\gesperrt{Abel}'schen gehören, dass alle \gesperrt{Abel}'schen Gleichungen durch eine
Kette von cyklischen Gleichungen vom Primzahlgrad und diese
letzteren durch Radicale lösbar sind.

Wir ersetzen also die Frage nach der Lösbarkeit durch Radicale
durch die damit gleichbedeutende der Lösbarkeit durch Wurzeln
\gesperrt{Abel}'scher Gleichungen. Ersetzen wir die \gesperrt{Abel}'sche Gleichung durch
eine Kette von cyklischen Gleichungen von Primzahlgrad, so wird
die erste Reduction der Gruppe durch Adjunction der Wurzel einer
solchen Gleichung eintreten, und wir stehen also zunächst vor der
Frage:

\gesperrt{Unter welchen Bedingungen wird die Gruppe $P$
einer Gleichung $n^{\mathrm{ten}}$ Grades $f(x) = 0$ durch Adjunction
einer Wurzel einer cyklischen Gleichung von Primzahlgrad $m$
reducirt?}

Wir beschränken uns hierbei nicht auf irreducible Gleichungen
$f(x) = 0$, sondern erörtern die Frage allgemein, immer unter der
selbstverständlichen Voraussetzung, dass $f(x)$ keine mehrfachen
Wurzeln hat.

Es ist denkbar, dass es nöthig ist, dem Körper $\Omega$ zunächst
verschiedene Wurzeln cyklischer Gleichungen (z.\,B. Einheitswurzeln)
zu adjungiren, die die Gruppe $P$ nicht verändern. Da es sich aber
jetzt nur um die Ermittelung der nothwendigen Eigenschaften der
Gruppe $P$ handelt, so nehmen wir an, es seien alle nöthigen
Vorbereitungen getroffen und der Körper $\Omega$ so beschaffen,
dass durch Adjunction einer Wurzel der cyklischen Gleichung
$\varphi(x) = 0$, von Primzahlgrad $m$, die Gruppe $P$ sich auf
einen ihrer Theiler, $Q$, reducirt.

Bezeichnen wir die Wurzeln von $\varphi(x) = 0$ mit

\[
\varepsilon,\ \varepsilon_1,\ \varepsilon_2 \ldots \varepsilon_{m-1},
\]
so sind alle diese Grössen rational (in $\Omega$) durch eine
beliebige unter ihnen ausdrückbar, und wenn $P$ die Gruppe von
$f(x) = 0$
%CONT>>
